\documentclass{article}
\usepackage[paperwidth=500mm,paperheight=700mm,margin=18mm]{geometry}
\usepackage[T1]{fontenc}
\usepackage{lmodern}
\usepackage{amsmath,amssymb}
% Scale large mathematical operators with the poster's math sizes.
\input{omxlmex.fd}
\DeclareFontShape{OMX}{lmex}{m}{n}{<->lmex10}{}
\usepackage{xcolor,colortbl}
\usepackage{tabularx,array}
\pagestyle{empty}
\definecolor{navy}{HTML}{172C3C}
\definecolor{coral}{HTML}{BF604D}
\definecolor{rulecolor}{HTML}{A5AFB6}
\color{navy}
\setlength{\parindent}{0pt}
\setlength{\parskip}{0pt}
\setlength{\arrayrulewidth}{0.6pt}
% Central typography controls, explicitly applied in table cells.
\newcommand{\TitleFont}{\fontsize{70}{76}\selectfont\bfseries}
\newcommand{\KickerFont}{\fontsize{28}{34}\selectfont\bfseries}
\newcommand{\HeadingFont}{\fontsize{36}{42}\selectfont\bfseries}
\newcommand{\DefinitionFont}{\fontsize{26}{30}\selectfont}
\newcommand{\SymbolFont}{\fontsize{28}{32}\selectfont}
\newcommand{\ExampleFont}{\fontsize{25}{30}\selectfont}
\newcommand{\TableHeadFont}{\fontsize{24}{29}\selectfont\bfseries}
\newcommand{\FooterFont}{\fontsize{24}{29}\selectfont}
\DeclareMathSizes{28}{28}{20}{17}
\DeclareMathSizes{26}{26}{19}{16}
\DeclareMathSizes{25}{25}{18}{15}

\newlength{\SymbolWidth}
\newlength{\MeaningWidth}
\newlength{\RowPadding}
\setlength{\RowPadding}{1.4mm}
% Examples for mu and sigma use the population data set {3,7}.
% Explicit line breaks keep related examples together inside a cell.
\newcommand{\exlines}[1]{\begin{tabular}[t]{@{}l@{}}#1\end{tabular}}
\newcommand{\symrow}[3]{%
  {\SymbolFont\rule[-3mm]{0pt}{11mm}\ensuremath{\displaystyle #1}}&%
  {\DefinitionFont #2}&{\ExampleFont #3}\\[\RowPadding]\hline
}
\newcommand{\sectionbox}[2]{%
  {\HeadingFont\color{coral}#1\par}%
  \vspace{2mm}{\color{navy}\hrule height 1.2pt}\vspace{2mm}%
  \begingroup
  \setlength{\SymbolWidth}{.23\linewidth}%
  \setlength{\MeaningWidth}{.36\linewidth}%
  \setlength{\tabcolsep}{2mm}%
  \renewcommand{\arraystretch}{1}%
  \arrayrulecolor{rulecolor}%
  \begin{tabularx}{\linewidth}{@{}>{\centering\arraybackslash\SymbolFont}p{\SymbolWidth}>{\raggedright\arraybackslash\DefinitionFont}p{\MeaningWidth}>{\raggedright\arraybackslash\ExampleFont}X@{}}
  {\TableHeadFont Symbol}&{\TableHeadFont Meaning}&{\TableHeadFont Example}\\[2mm]\hline
  #2
  \end{tabularx}%
  \endgroup\par\vspace{5mm}%
}
\begin{document}
{\KickerFont THE LANGUAGE OF\par}
\vspace{1mm}
{\TitleFont\color{coral}MATHEMATICS\par}
\vspace{5mm}
{\color{navy}\hrule height 2pt}
\vspace{8mm}
\noindent
\begin{minipage}[t]{226mm}
\sectionbox{General notation}{
\symrow{=}{Equal to}{$2+3=5$}
\symrow{\neq}{Not equal to}{$2\neq 3$}
\symrow{\approx}{Approximately equal to}{$\pi\approx 3.14$}
\symrow{<,\ >}{Less than; greater than}{$2<3;\quad 5>4$}
\symrow{\leq,\ \geq}{Less than or equal to; greater than or equal to}{\exlines{$2\leq 2$\\$5\geq 4$}}
\symrow{\pm}{Plus or minus}{$3\pm 1=4\text{ or }2$}
\symrow{|a|}{Absolute value}{$|-5|=5$}
\symrow{\sqrt{a}}{Square root}{$\sqrt{9}=3$}
}
\sectionbox{Sets and number systems}{
\symrow{\{a,b,c\}}{Set with three elements}{$\{2,4,6\}$}
\symrow{\varnothing}{Empty set}{$\{1\}\cap\{2\}=\varnothing$}
\symrow{\in,\ \notin}{Is an element of; is not an element of}{\exlines{$2\in\{1,2\}$\\$3\notin\{1,2\}$}}
\symrow{\subseteq}{Is a subset of}{$\{1\}\subseteq\{1,2\}$}
\symrow{\cup,\ \cap}{Union; intersection}{\exlines{$\{1\}\cup\{2\}=\{1,2\}$\\$\{1,2\}\cap\{2\}=\{2\}$}}
\symrow{\mathbb N,\ \mathbb Z}{Natural numbers; integers}{$3\in\mathbb N;\quad -2\in\mathbb Z$}
\symrow{\mathbb Q,\ \mathbb R}{Rational numbers; real numbers}{$\frac12\in\mathbb Q;\quad\sqrt{2}\in\mathbb R$}
\symrow{(a,b),\ [a,b]}{Open interval; closed interval}{\exlines{$1\notin(1,3)$\\$1\in[1,3]$}}
}
\sectionbox{Logic}{
\symrow{\Rightarrow}{Implies}{$x>2\Rightarrow x>1$}
\symrow{\Leftrightarrow}{If and only if}{$x+1=3\Leftrightarrow x=2$}
\symrow{\forall}{For every}{$\forall x\in\mathbb R,\ x^2\geq0$}
\symrow{\exists}{There exists}{$\exists x\in\mathbb Z:\ x+2=0$}
\symrow{\neg}{Logical negation}{$\neg(2>3)$ is true}
}
\sectionbox{Summation notation}{
\symrow{\sum_{i=1}^{n}a_i}{Add $a_i$ for each integer $i$ from $1$ to $n$.}{$\sum_{i=1}^{3}2i=2+4+6$}
}
{\DefinitionFont
In $\sum_{i=1}^{n}a_i$, $i$ is the \textbf{index}, $1$ is the
\textbf{lower bound}, $n$ is the \textbf{upper bound}, and $a_i$ is
 the \textbf{expression being summed}.\par
\vspace{3mm}
{\color{coral}\bfseries Worked examples}\par
\vspace{2mm}
$\displaystyle\sum_{i=1}^{4}i=1+2+3+4=10$\par
\vspace{3mm}
$\displaystyle\sum_{k=1}^{3}k^2=1+4+9=14$\par
}
\end{minipage}\hfill
\begin{minipage}[t]{226mm}
\sectionbox{Probability and events}{
\symrow{\Omega}{Sample space: all possible outcomes}{$\Omega=\{1,2,3,4,5,6\}$}
\symrow{\omega}{An individual outcome}{$\omega=3\in\Omega$}
\symrow{A=\{1,3,5\}}{Event shown as a set of outcomes}{$A=\{1,3,5\}$: odd roll}
\symrow{A\subseteq\Omega}{An event is a subset of the sample space}{$\{1,3,5\}\subseteq\Omega$}
\symrow{P(A)}{Probability of event $A$}{$P(A)=\frac36=\frac12$}
\symrow{A^c}{Complement of event $A$}{$A^c=\{2,4,6\}$}
\symrow{A\cup B}{Event $A$ or $B$ (or both)}{\exlines{$B=\{1,2\}$\\$A\cup B=\{1,2,3,5\}$}}
\symrow{A\cap B}{Both events $A$ and $B$}{$A\cap B=\{1\}$}
\symrow{P(A\mid B)}{Probability of $A$ given $B$}{$P(A\mid B)=\frac12$}
}
{\ExampleFont Probability examples use a fair six-sided die,\par
$A=\{1,3,5\}$ and $B=\{1,2\}$.\par}
\vspace{5mm}
\sectionbox{Function notation}{
\symrow{f}{Name of a function}{$f(x)=x^2$}
\symrow{f(x)}{Value of $f$ at input $x$}{$f(3)=3^2=9$}
\symrow{f:A\to B}{Function from set $A$ to set $B$}{$f:\mathbb R\to\mathbb R,\ f(x)=x^2$}
\symrow{x\mapsto f(x)}{Input-to-output mapping}{$x\mapsto x^2;\quad 3\mapsto9$}
\symrow{\operatorname{dom}(f)}{Domain of a function}{$\operatorname{dom}(f)=\mathbb R$}
}
\sectionbox{Letters and Greek symbols}{
\symrow{a,b,c,\ldots}{Common lowercase Latin letters}{$a=2,\ b=3,\ c=a+b$}
\symrow{A,B,C,\ldots}{Common uppercase Latin letters}{\exlines{$A=\{1\},\ B=\{2\}$\\$C=A\cup B$}}
\symrow{x,y,z}{Other common Latin letters}{$y=2x,\quad z=x+y$}
\symrow{\alpha,\ \beta,\ \gamma}{Alpha; beta; gamma}{\exlines{Triangle angles:\\$\alpha+\beta+\gamma=180^\circ$}}
\symrow{\delta,\ \Delta}{Delta; uppercase delta}{$\delta=0.1;\quad\Delta x=x_2-x_1$}
\symrow{\theta,\ \Theta}{Theta; uppercase theta}{$\theta=30^\circ;\quad\Theta=2\theta$}
\symrow{\lambda,\ \Lambda}{Lambda; uppercase lambda}{$\lambda=2;\quad\Lambda=\{1,2\}$}
\symrow{\mu,\ \sigma,\ \Sigma}{Mu; sigma; uppercase sigma}{\exlines{$\mu=\frac{3+7}{2}=5$\\$\sigma=\sqrt{4}=2$\\$\sum_{i=1}^{3}i=6$}}
\symrow{\pi,\ \Pi}{Pi; uppercase pi}{\exlines{$\pi\approx3.14$\\$\prod_{j=1}^{3}j=6$}}
\symrow{\phi,\ \Phi}{Phi; uppercase phi}{$\phi=\frac{1+\sqrt5}{2};\quad\Phi=\phi^2$}
\symrow{\omega,\ \Omega}{Omega; uppercase omega}{$\Omega=\{H,T\};\quad\omega=H$}
}
\end{minipage}
\vfill
{\FooterFont Symbols acquire meaning from their mathematical context.\par}
\end{document}



